\documentclass[acmsmall,screen,review,anonymous,nonacm]{acmart}

\usepackage{amsmath}
\usepackage{booktabs}
\usepackage{tikz}
\usetikzlibrary{fit,tikzmark}
\usepackage{xcolor}
\usepackage{xspace}
\usepackage[ruled,vlined,linesnumbered]{algorithm2e}
\usepackage{wrapfig}
\usepackage{needspace}
\usepackage{svg}
% Draft-only spacing overrides; remove for ACM submission.
\setlength{\textfloatsep}{10pt plus 2pt minus 2pt}
\setlength{\floatsep}{8pt plus 2pt minus 2pt}
\setlength{\intextsep}{10pt plus 2pt minus 2pt}
\definecolor{klensblue}{HTML}{1764AB}
\definecolor{ksearchgray}{HTML}{6B7280}
\definecolor{kernelbandorange}{HTML}{E07A1F}
\definecolor{openevolvegreen}{HTML}{2A9D6F}
\definecolor{codexpurple}{HTML}{8B5FBF}

\newcommand{\klens}{\textsc{K-LENS}\xspace}
\newcommand{\ksearch}{\textsc{K-Search}\xspace}
\newcommand{\kernelband}{\textsc{KernelBand}\xspace}
\newcommand{\openevolve}{\textsc{OpenEvolve}\xspace}
\newcommand{\standardcodex}{\textsc{Codex}\xspace}
\newcommand{\ascendc}{\textsc{AscendC}\xspace}
\newcommand{\nolessons}{\textsc{No-Lessons}\xspace}
\newcommand{\onlycodinglessons}{\textsc{Only-Coding}\xspace}
\newcommand{\tbd}[1]{\textcolor{red}{[TBD: #1]}}
\newcommand{\numkernels}{16\xspace}
\newcommand{\numprofiles}{100\xspace}
\newcommand{\numreps}{3\xspace}
\newcommand{\numruns}{192\xspace}
\newcommand{\numtotalprofiles}{15{,}360\xspace}
\newcommand{\dimitris}[1]{{\color{blue}\textbf{(Dimitris:} #1\textbf{)}}}

% TAPS-safe summary box, matching the result boxes in gordian.tex.
\newsavebox{\klensresultcontent}
\newenvironment{klensresultbox}{%
  \par\addvspace{0.5\baselineskip plus 2pt}%
  \noindent
  \begingroup
  \setlength{\fboxrule}{0.3mm}%
  \setlength{\fboxsep}{\dimexpr1mm+2pt\relax}%
  \begin{lrbox}{\klensresultcontent}%
  \begin{minipage}{\dimexpr\linewidth-2\fboxrule-2\fboxsep\relax}%
}{%
  \end{minipage}%
  \end{lrbox}%
  \fcolorbox{black!20}{gray!5}{\usebox{\klensresultcontent}}%
  \endgroup
  \par\addvspace{0.5\baselineskip plus 2pt}%
}

\title{K-LENS: Evidence-Reranked Lessons for Accelerator Kernel Optimization}


\author{Dimitrios Stamatios Bouras}
\orcid{0009-0001-3171-5194}
\email{2501112125@stu.pku.edu.cn}
\affiliation{%
  \department{Key Lab of HCST (PKU), MoE and School of Computer Science}
  \institution{Peking University}
  \city{Beijing}
  \country{China}
}
\affiliation{%
  \department{Global Technical Service (GTS)}
  \institution{Huawei Technologies Co., Ltd}
  \city{Shanghai}
  \country{China}
}

\author{Jian Meng}
\orcid{0000-0002-7703-5020}
\email{mengjian11@huawei.com}
\affiliation{%
  \company{Huawei Technologies}
  \city{Shanghai}
  \country{China}}

\author{Junjie Huang}
\orcid{0009-0007-7369-1788}
\email{huangjunjie76@huawei.com}
\affiliation{%
  \institution{Huawei Technologies Co., Ltd.}
  \city{Shenzhen}
  \country{China}}

\author{Jing Li}
\orcid{0000-0001-9933-8383}
\email{lijing473@huawei.com}
\affiliation{%
  \institution{Huawei Technologies Co., Ltd.}
  \city{Shanghai}
  \country{China}}


\author{Zhanglian Lin}
\orcid{0009-0004-2303-5211}
\email{linzhanglian@huawei.com}
\affiliation{%
  \institution{Huawei Technologies Co., Ltd.}
  \city{Shanghai}
  \country{China}
}

\author{Zhicheng Zhou}
\orcid{0009-0005-4431-0611}
\email{zhouzhicheng15@huawei.com}
\affiliation{%
  \institution{Huawei Technologies Co., Ltd.}
  \city{Shenzhen}
  \country{China}}

\author{Gengpei Qi}
\orcid{0009-0002-2086-6859}
\email{qigengpei@huawei.com}
\affiliation{%
  \department{}
  \institution{Huawei Technologies Co., Ltd.}
  \city{Shanghai}
  \country{China}}

\author{Linyuan Qiu}
\orcid{0009-0003-1899-3063}
\email{qiulinyuan@huawei.com}
\affiliation{%
  \institution{Huawei Technologies Co., Ltd.}
  \city{Shenzhen}
  \country{China}}

\author{Abdullin Amir}
\orcid{0009-0006-0889-2973}
\email{abdullin.amir@h-partners.com}
\affiliation{%
  \company{Huawei Technologies Co., Ltd.}
  \city{Shanghai}
  \country{China}}

\author{Ivan Nosov}
\orcid{0009-0003-7231-6274}
\email{nosov.ivan@h-partners.com}
\affiliation{%
  \company{Huawei Technologies Co., Ltd.}
  \city{Shanghai}
  \country{China}}

\author{Bekzat Rakhimbayev}
\orcid{0009-0005-0216-4951}
\email{rakhimbayev.bekzat@huawei.com}
\affiliation{%
  \company{Huawei Technologies Co., Ltd.}
  \city{Shanghai}
  \country{China}}

\author{Sergey Mechtaev}
\authornote{Sergey Mechtaev is the corresponding author.}
\orcid{0000-0001-6088-4993}
\email{mechtaev@pku.edu.cn}
\affiliation{%
  \department{Key Lab of HCST (PKU), MoE and School of Computer Science}
  \institution{Peking University}
  \city{Beijing}
  \country{China}
}
\renewcommand{\shortauthors}{Bouras et al.}

\begin{document}


\begin{abstract}
LLM-guided optimization of low-resource domain-specific languages, such as
\ascendc for Ascend NPU kernels, is hard: models are unfamiliar with the
language's performance idioms, and each evaluation requires compiling,
testing, and profiling on the device. A search can therefore afford few
evaluations and gets little guidance from the model, so every measurement
must become reusable knowledge rather than a one-off verdict.
We present \klens, a search system that supplements the profiled kernel
implementations it explores with measured optimization lessons. \klens extracts a lesson from each evaluated optimization, recording
the conditions of the optimized candidate, the strategy applied, and the
measured outcome. For each subsequent candidate, \klens generates optimization
proposals using the lessons whose recorded conditions match it, and ranks them
by the accumulated performance evidence of the lessons behind them; each
measured result strengthens or weakens this evidence and reranks proposals.
To initialize a new search, the ``warm-start'' version of \klens imports
lessons extracted from similar candidates in previous searches on other
kernels.
We evaluate \klens on \numkernels production \ascendc kernels spanning
attention, normalization, and mixture-of-experts operators, whose original
implementations already contain hardware-specific tuning. We compare against
\ksearch, \kernelband, \openevolve, and the \standardcodex coding harness under
equal candidate-evaluation and LLM-token budgets.
With \numprofiles candidate evaluations per kernel, \klens achieves a
geometric-mean speedup of 22.40\% over the original implementations, nearly
double the strongest baseline's 11.38\%, and the highest speedup at every
reported token budget.
Removing the lesson layer reduces the gain to 13.15\%, indicating that lessons
account for most of \klens's advantage over the baselines.
With warm start, the gain at 20 candidate evaluations rises from 17.58\% to
28.45\%.
These results suggest that accumulating and transferring measured
knowledge enables effective LLM-guided optimization in low-resource languages.
\end{abstract}

\begin{CCSXML}
<ccs2012>
  <concept>
    <concept_id>10011007.10011074.10011099</concept_id>
    <concept_desc>Software and its engineering~Software performance</concept_desc>
    <concept_significance>500</concept_significance>
  </concept>
  <concept>
    <concept_id>10010147.10010257</concept_id>
    <concept_desc>Computing methodologies~Machine learning</concept_desc>
    <concept_significance>500</concept_significance>
  </concept>
</ccs2012>
\end{CCSXML}

\ccsdesc[500]{Software and its engineering~Software performance}
\ccsdesc[500]{Computing methodologies~Machine learning}
\keywords{kernel optimization, large language models, search-based software engineering, AscendC}

\maketitle
\section{Introduction}

LLM-guided code optimization relies on knowledge of a language's performance
idioms and feedback from executing candidate programs. Specialized
domain-specific languages can offer little of either: public examples are
scarce, and evaluating a transformation may require substantial implementation
work. Under a limited budget, choosing a promising strategy is critical.

Accelerator kernels implement computations on devices such as GPUs and neural
processing units (NPUs). Optimizing these programs requires coordinated edits
across multiple files, followed by compilation, testing, and profiling on the
device. These requirements are especially demanding in complex kernel
languages with scarce public code examples, such as \ascendc, Huawei's
C++-based kernel language for Ascend NPUs~\citep{Cao2026AscendKernelGen}.
MultiKernelBench shows this difficulty: under one-shot prompting,
evaluated LLMs achieve correctness rates (Pass@1) of 0.7--2.5\% on \ascendc,
compared with 15.8--52.6\% on CUDA~\citep{Wen2025MultiKernelBench}.
Optimization inherits this difficulty: a candidate must compile and pass tests
before its latency is measured, and correct kernels often fail to beat a
baseline~\citep{Ouyang2025KernelBench}.

A developer or an LLM can suggest optimization strategies, high-level code
changes such as moving loop-invariant computations outside a loop, at
relatively low cost. Implementing each strategy as working kernel code and
evaluating it, however, can require substantial time, human effort, or LLM
tokens. Under a fixed evaluation budget, the search's effectiveness therefore
depends on carefully selecting which strategies to pursue. In a
low-resource language, the LLM's own judgment offers limited guidance for this
selection, as it reflects the same unfamiliarity with the language's
performance idioms that makes these kernels difficult to optimize.

Prior search-based optimizers reuse measurements to guide search, but only
within the transformations and parameters encoded in their search
spaces~\citep{Adams2019Halide,Zheng2020Ansor,Ansel2014OpenTuner}; they do not
extract explicit information about which code changes helped under which
conditions. More recent approaches combine LLMs' ability to propose
source-code changes with guidance from earlier attempts, stored as
optimization histories, reusable lessons, or
skills~\citep{Liu2025Lessons,Dong2026KernelBlaster,Sun2026KernelSkill}. Such
guidance must account for differences between implementations: a strategy that
improves one may reduce performance in another with different bottlenecks.
Identifying relevant guidance, however, does not determine which proposal
deserves implementation and profiling effort. When an LLM makes this decision
by reviewing previous attempts, every new measurement requires reinterpreting
earlier proposals, so inference cost grows with the search. For example,
\ksearch asks an LLM to assign priorities to proposals and revises them after
each cycle, reasoning over a growing
context~\citep{Cao2026K}. Two challenges remain: \emph{lesson reuse},
applying a measurement on one candidate to proposals on others, and the
\emph{cost of re-reasoning}, updating priorities without the LLM reassessing
the whole search state.

We present \klens (Kernel Lesson-ENhanced Search), a search system that uses
measured results to propose and rank optimization strategies and revise their
priorities as new results become available. \klens extracts reusable
\emph{strategy lessons} from optimization results. Each lesson records the
characteristics of the candidate being optimized, the strategy attempted, and
the observed outcome. These characteristics help identify relevant lessons
when proposing an optimization strategy for another candidate. Each proposed
optimization strategy ``cites'' the lessons that informed it. Each lesson has
a \emph{reputation} score that accumulates signed performance changes relative
to the candidate being modified. \klens uses these scores to rank proposed
optimization strategies awaiting implementation. New measurements update the
cited lessons' reputations and the rankings of other proposed strategies that
cite them. A result from one candidate can therefore redirect optimization
effort elsewhere in the search. These numerical updates avoid regenerating
proposals or repeatedly asking the LLM to reassess them and the search context
as the search grows. Complementary \emph{coding lessons}, extracted from
diagnostics and code changes across attempts to implement the same strategy, record repairs and
constraints that help later attempts implement a chosen strategy correctly.
Strategy lessons guide what to try; coding lessons guide how to implement it.

Lessons record their measurement conditions, so they outlive the search and
yield reusable knowledge, not just a faster kernel. \klens can therefore start a new search with lessons from other
kernels, a \emph{warm start}, which matters most under small budgets, selecting
them by kernel and bottleneck similarity, measured gains, and reputation.

We evaluate \klens using GPT-5.6 Luna on 16 production \ascendc kernels, with
implementations spanning up to 35 files and 10{,}382 SLOC, against three
search-based optimization methods (\ksearch, \kernelband, and \openevolve) and
the general-purpose \standardcodex coding agent. At 100 evaluation attempts,
\klens achieves a 22.40\% gain in geometric-mean speedup, nearly twice the
strongest baseline's 11.38\%, and leads at every reported LLM-token budget.
Removing the lesson layer reduces the gain to 13.15\%. At 20 evaluation
attempts, transferring lessons from other kernels increases the gain from
17.58\% to 28.45\%, exceeding the geometric-mean gain without transfer at 100
attempts. A comparison using Kimi K3 on six complex kernels, those exceeding
3{,}000 SLOC, also favors \klens, providing evidence that its advantage extends to another model
family.

This paper contributes:
\vspace{-3pt}
\begin{itemize}
  \item \textbf{A lesson representation} for optimizing programs in low-resource
  languages like \ascendc. These strategy lessons record the conditions of the
  candidate being optimized, the strategy applied, and its measured outcome.
  \item \textbf{A search method} that uses lessons to propose and rank optimization
  strategies, then uses measured performance changes to update lesson scores
  and rerank pending proposals.
  \item \textbf{An evaluation} on \ascendc kernels against three search approaches and
  the Codex harness under equal evaluation-attempt and
  inference budgets, including a lesson-layer ablation, transfer of lessons
  across kernels, and repeated runs to assess LLM variability.
\end{itemize}








% \section{Introduction}

% LLM-guided code optimization relies on knowledge of a language's performance idioms and feedback from
% executing candidate programs. Specialized domain-specific languages can offer little of either:
% public examples are scarce, and evaluating a transformation may require substantial implementation
% work. Under a limited budget, choosing a promising strategy is critical.

% Accelerator kernels implement computations on devices such as GPUs and
% neural processing units (NPUs). Optimizing these programs requires
% coordinated edits across multiple files, followed by compilation,
% testing, and profiling on the device. These requirements can be
% particularly demanding in complex kernel languages with scarce public
% code examples, such as \ascendc, Huawei's C++-based kernel language for
% Ascend NPUs~\citep{Cao2026AscendKernelGen}.
% MultiKernelBench illustrates this difficulty: under one-shot prompting,
%   the evaluated LLMs achieve correctness rates (Pass@1) of 0.7--2.5\%
%   on \ascendc, compared with 15.8--52.6\% on
%   CUDA~\citep{Wen2025MultiKernelBench}. A developer or an LLM can suggest optimization strategies, high-level
% code changes such as moving loop-invariant computations outside a
% loop, at relatively low cost. Implementing each strategy as working
% kernel code and evaluating it though, can require substantial time, human effort, or LLM tokens.
% The quality of the search for optimised kernel implementations therefore
% depends on carefully selecting which strategies to pursue.


% Earlier search-based code optimization methods use measured or predicted
% performance to iteratively improve program implementations. Tree-based
% methods explore sequences of program
% transformations~\citep{Adams2019Halide}, population-based methods retain
% and modify sets of implementations across
% generations~\citep{Zheng2020Ansor}, and bandit-based methods allocate
% evaluations among search techniques~\citep{Ansel2014OpenTuner}.
% These methods reuse measurements to guide search, but their choices
% remain limited to the transformations and parameters encoded in their
% search spaces. They do not extract explicit, reusable lessons about
% which code changes improved performance under which conditions.
% More recent approaches extend these search methods by combining LLMs'
% ability to propose source-code changes with guidance from earlier
% attempts, stored as optimization histories, reusable lessons, or
% skills~\citep{Liu2025Lessons,Dong2026KernelBlaster,Sun2026KernelSkill}.
% This guidance must account for differences between implementations:
% a strategy that improves one may reduce performance in another with
% different bottlenecks. However, identifying relevant guidance does not
% itself determine which proposed optimization should receive implementation
% and profiling effort. When this decision relies on an LLM reviewing
% previous attempts, new measurements must be interpreted again to
% reconsider earlier proposals. As the search grows, this requires repeated
% LLM reasoning over more results and alternatives, which increases
% inference cost, limits scalability, and may degrade the
% proposal selection quality.
% For example, \ksearch uses an LLM to assign priorities to proposals and based on the search state and execution feedback revises these priorities after each cycle, 
% requiring repeated reasoning over a growing search and proposal context~\citep{Cao2026K}.
% The remaining challenge is to use measured outcomes to update proposal
% priorities directly, without regenerating proposals or requiring repeated
% LLM reasoning over the whole search state.

% % We present \klens (Kernel Lesson-ENhanced Search), an iterative multi-agent
% %   search system that ranks optimization proposals using reusable
% %   \emph{strategy lessons} and updates those rankings as new performance
% %   measurements become available. The search maintains profiled kernel
% %   implementations, called \emph{candidates}. A proposal specifies a
% %   high-level optimization strategy for one candidate. The search proceeds
% %   in cycles: each cycle selects one proposal for implementation and
% %   evaluates the resulting code to produce a new candidate. Each strategy
% %   lesson records the conditions of the candidate being optimized, the
% %   strategy attempted, and the observed outcome. Each lesson also stores a
% %   \emph{reputation} score that accumulates signed performance changes,
% %   each measured relative to the candidate being modified. At the start
% %   of a cycle, a planner subagent selects promising candidates
% %   and generates proposals using relevant lessons identified by comparing
% %   the candidates' conditions and profiling results with those recorded
% %   in the available lessons. Each proposal
% %   identifies the lessons that influenced its strategy; we call these
% %   references \emph{lesson citations}.

% %   Proposals awaiting implementation and evaluation remain pending.
% %   \klens ranks pending proposals using the reputation scores of cited
% %   lessons. It deterministically forms
% %   a shortlist that prioritizes highly ranked proposals while
% %   preserving candidate diversity, limiting how many proposals the selector
% %   considers as the search grows. A selector subagent compares the shortlisted
% %   proposals and selects the one it considers most promising for improving
% %   performance. For the selected proposal, an implementor subagent makes a
% %   bounded number of implementation attempts. \klens compares the fastest
% %   correct, profiled implementation with the original candidate to compute
% %   the performance change. It uses this change to update the reputations
% %   of the cited lessons used to rank that proposal and extracts a new
% %   strategy lesson from the new candidate.
% %   These updates rerank pending proposals that cite the updated lessons.
% %   A measured result from one candidate can therefore change which strategy
% %   is selected for another, without regenerating proposals or asking an LLM
% %   to reassess each one.
% We present \klens (Kernel Lesson-ENhanced Search), a search system that
%   uses measured results to propose and rank optimization strategies and
%   revise their priorities as new results become available.
%   \klens extracts
%   reusable \emph{strategy lessons} from optimization results, wich record the characteristics of the kernel candidate being optimized,
%   the strategy attempted, and the observed outcome. These characteristics
%   help identify relevant lessons when proposing an optimization strategy
%   for another candidate. Each proposed optimization strategy ``cites''
%   the lessons that informed it, explicitly recording which lessons
%   influenced the choice of strategy. Each lesson has a
%   \emph{reputation} score that accumulates signed performance changes
%   relative to the candidate being modified. \klens uses these scores
%   to rank proposed optimization strategies awaiting implementation.
%   New measurements update the cited lessons' reputations and the rankings
%   of other proposed strategies that cite them. A result from one candidatecan therefore redirect optimization effort elsewhere in the search.
%   These numerical updates avoid regenerating proposals or repeatedly asking
%   the LLM to reassess them and the search context as the search grows.

%   The same lesson representation also supports reuse across different kernels.
%   \klens can initialize a new search with lessons extracted from
%   kernel candidates in previous searches on new kernels, a process we call
%   a \emph{warm start}. It selects lessons using similarities in kernel
%   characteristics and performance bottlenecks, measured performance gains,
%   and supported lesson reputation. Measurements from the new search update the
%   transferred lessons' reputations as the search produces lessons of its own.

% % To address this problem , We present \klens (Kernel Lesson-ENhanced Search), which extacts and stores previous
% %   optimization results as reusable \emph{strategy lessons}. Each lesson records
% %   the conditions of the kernel candidate being modified, the optimization strategy
% %   applied, and the observed outcome. Each lesson also stores an
% %   \emph{implementation difficulty} score based on implementation attempts and
% %   evaluation time. The lesson also stores a \emph{reputation}
% %   score that accumulates signed performance changes measured relative to the
% %   candidates being modified. A planner subagent uses these lessons and current profiled candidates to generate proposals for next optimisations. \dimitris{what are next optimizations?} A proposal specifies an optimization strategy for an existing candidate \dimitris{what is an existing candidate?} and "cites" the lesson(s) that influenced that choice of strategy.
% %   \klens uses the cited lessons to rank pending proposals. \dimitris{what are pending proposals?} Each proposal's rank
% %   depends on the reputation scores of the lessons it cites.
% %   It then forms a bounded shortlist \dimitris{can a shortlist be unbounded?} from the
% %   ranked proposals, adding proposals for promising candidates and candidates
% %   infrequently represented. A selector agent then chooses a proposal from this
% %   shortlist using the corresponding candidates' profiling results.
% %   \klens computes proposal scores numerically and uses fixed rules to form
% %   the shortlist. \dimitris{fixed rules sounds like ``magic''} The selector therefore compares a bounded number of proposals as the search grows.

% %   For the selected proposal, an implementor agent makes a bounded number of implementation attempts and compares the fastest correct implementation
% %   with the original candidate to compute the performance change. \klens then
% %   uses this change to update the reputation scores of the
% %   lessons cited by this proposal and extracts a new strategy lesson describing the conditions, strategy,
% %   and outcome. These updates change the rankings of pending proposals that also cite
% %   the updated lessons without requiring the planner to generate those proposals
% %   again. The shortlist limits how many pending proposals the selector considers as
% %   thier number grows, while including alternatives for unrepresented candidates
% %   encourages broader exploration.
% % If all implementation attempts fail, \klens
% % retains diagnostics and code changes for extracting coding lessons but assigns
% % no performance penalty to the strategy because its performance was not measured.

% % Complementary \emph{coding lessons} turn diagnostics and code changes between implementation attempts
% % for the same proposal into reusable repairs and implementation constraints. They guide the implementor subagent in realizing
% % a chosen strategy, preserving useful experience even when no correct candidate is produced. Strategy
% % lessons guide what to try; coding lessons guide how to implement it and do not enter proposal ranking or selection. \dimitris{overall, these paragraphs look like a method description rather than introduction. an introduction needs to focus on key ideas and innovations, and how they enabled improvements over previous work.}

% % Since each strategy lesson is conditioned on the properties of the
% % candidate whose optimization produced that lesson, \klens can retrieve
% % relevant lessons for proposal generation by comparing those properties
% % with those of the new candidate. This comparison also allows \klens to
% % initialize a new search with lessons from similar candidates in previous
% % searches on other kernels, a process we call a \emph{warm start}.
% % To select a fixed-size set of lessons for transfer,
% % \klens considers candidate similarity, measured
% % performance gains, and lesson reputation. Measurements from the new
% % search update the transferred lessons' reputation scores, while the
% % search also produces its own lessons.

% We evaluate \klens using GPT-5.6 Luna on 16 production \ascendc kernels,
% with implementations spanning up to 35 files and 10{,}382 SLOC,
% against three search-based optimization methods (\ksearch, \kernelband,
% and \openevolve) and the general-purpose \standardcodex coding agent.
% At 100 evaluation attempts, \klens achieves a 22.40\% gain in
% geometric-mean speedup, nearly twice the strongest baseline's 11.38\%,
% and leads at every reported LLM-token budget. Removing the lesson layer
% reduces the gain to 13.15\%. At 20 evaluation attempts, transferring lessons from
% other kernels increases the gain from 17.58\% to 28.45\%.
% A comparison with Kimi K3 on seven of the largest kernels also favors
% \klens, providing evidence that its advantage extends to another model family.

% This paper contributes:
% \vspace{-3pt}
% \begin{itemize}
%  \item A lesson representation for optimizing programs in low resource languages like \ascendc. These strategy lessons record the conditions of the candidate being
%   optimized, the strategy applied, and its measured outcome.
%   \item A search method that uses lessons to propose and rank optimization
%   strategies, then uses measured performance changes to update lesson scores
%   and rerank pending proposals.
%   \item An evaluation on \ascendc kernels against three search approaches
%   and a general-purpose Codex harness under equal candidate-evaluation and
%   inference budgets, including a lesson-layer ablation, transfer of lessons
%   across kernels, and repeated runs to assess LLM variability.
% \end{itemize}





\section{Motivating Example}
\label{sec:motivation}
An \emph{operator} defines a computation on multidimensional arrays, called
tensors; a \emph{kernel} is a hardware-specific program that implements it.
Our running example is the kernel of the \texttt{SparseAttnSharedkv} operator,
which fuses sparse attention (scoring selected query--key pairs, normalizing
the scores, and combining the corresponding values) into one kernel.
Figure~\ref{fig:motivating-example} illustrates how \klens extracts lessons from measured results and reuses them to guide the search toward better candidates.

\begin{figure}[t]
    \centering
    \includegraphics[width=\textwidth]{figures/motivating_example/k_lens_motivating_example.pdf}
    \caption{Reusing a strategy lesson in \texttt{SparseAttnSharedkv}. Code is simplified pseudocode.}
    % Without supporting evidence, proposal P10 is selected only after other
    % proposals fail, and it needs four implementation attempts; its 3.29\%
    % latency reduction is recorded as lesson L11. On candidate B, which has a
    % similar bottleneck and different optimizations, L11's reputation ranks
    % proposal P40 first, and every implementation attempt is correct. This yields a
    % further 7.02\% reduction (27.99\% from the baseline) and raises L11's
    % reputation. Code is simplified pseudocode.}
    \Description{A search tree, in which circles are candidates and arrows with
    a P box are proposals, links candidates A, A', B and B'. Candidate A
    re-checks unchanged conditions in every loop iteration; its child A' checks
    them once and runs a specialized loop, reducing latency by 3.29\%. Proposal
    P10, which produced A', had no supporting lesson: it was chosen only after
    proposals A1 and A2 failed, and its implementation failed three times before
    the fourth attempt was correct. The transition is recorded as strategy
    lesson L11 with reputation 0.03. Candidate B, on another branch, has the same
    bottleneck, so L11 creates proposal P40 on B. In the shortlist, which the
    LLM selector reads instead of the whole tree, P40 ranks first by its
    lesson's reputation, ahead of proposals P23 and P36 on faster parents, and
    is chosen immediately. Guided by prior coding lessons, all three attempts
    are correct; the best, B', is 7.02\% faster than B and 27.99\% faster than
    the baseline. The result raises L11's reputation from 0.03 to 0.10.}
    \label{fig:motivating-example}
\end{figure}

The kernel's hot path is a general loop that supports every input format the
operator accepts. Because it must handle any of them, it re-derives its
data-transfer plan and re-checks path and shape conditions on every iteration,
although these never change within a run. The optimization we want the search
to find specializes this loop: check those conditions once and run a fixed
schedule, while all other inputs keep the general loop.

At candidate A, an LLM proposed this specialization as proposal P10, one of
many proposals on A and on other candidates. Without a deep
understanding of \ascendc performance idioms, the LLM could only guess which
proposal was most promising, so the search first tried several others, such as
A1 and A2, which did not yield improvements. When P10 was finally selected,
no relevant coding lesson existed, and its implementation failed to
compile three times before the fourth attempt produced child A$'$, 3.29\%
faster than A. The lack of evidence therefore cost search cycles, and the lack of
implementation guidance then cost failed attempts.

\klens records this measured transition as strategy lesson L11 with the conditions it was applied under (A's
bottleneck and profile), the strategy, and its outcome. Its reputation, the sum of its measured parent-relative effects,
starts at 0.03; the compilation failures produce no measurement and only raise
the lesson's difficulty. It also records coding lessons on how to move the
checks out of the loop while maintaining correctness.

Candidate B lies on a different branch, retains different optimizations, and is
already 22.55\% faster than the baseline, but its profile shows the same
scalar-control bottleneck. \klens therefore retrieves L11, and the LLM creates
proposal P40 citing it. The ranker places P40 at the top of the shortlist, the
handful of open proposals that the LLM selector inspects instead of the whole
tree, because L11 carries a positive measured reputation (0.03). Competing
proposals such as P23 and P36 cite a weaker lesson or none (0.01 and 0.00),
although their parents are faster (3.647 and 3.513\,ms, against 3.762\,ms for
B). The selector picks P40 immediately.

Guided by prior coding lessons, such as checking every condition the
specialized loop assumes and keeping each buffer alive until its last use, all
three attempts on B compile and pass the tests, at 3.592, 3.533 and
3.498\,ms: the first is already correct, and the others extend the
specialization and repair its synchronization. The best, B$'$, is 7.02\%
faster than B and 27.99\% faster than the baseline.

Finally, \klens credits B's measured 7.02\% latency reduction to L11, raising
its reputation from 0.03 to 0.10, now from two measurements, so later proposals
that cite L11 rank higher still. The decision to implement P40 thus came from a
measured outcome on another candidate, not from an LLM re-reading the search
state: what took A many search cycles and four attempts took B one selection
and a first attempt that was already correct.

\section{K-LENS}
\label{sec:approach}

\subsection{Overview}

Figure~\ref{fig:overview} and Algorithm~\ref{alg:search} show the budgeted search
loop. \klens maintains a search tree of profiled kernel implementations,
called \emph{candidates}. It also stores pending \emph{proposals} and sets of
strategy and coding lessons. Each pending proposal specifies an optimization
strategy for a candidate; the proposal has not yet been implemented
or evaluated. The \emph{controller} manages these records and deterministically
updates proposal rankings and lesson reputations.

A planner agent begins each cycle by choosing candidates based on their measured
performance and opportunities for further improvement. For each candidate,
it uses lessons extracted from candidates with similar characteristics and
performance bottlenecks to generate proposals
(\texttt{Plan} in Algorithm~\ref{alg:search}).
Each proposal identifies the lessons that influenced the planner's choice of
optimization strategy. We call these references \emph{lesson citations}.
A citation indicates whether the proposal \emph{applies}
the lesson's strategy, \emph{adapts} it to the candidate, or uses the lesson to
\emph{avoid} an optimization. The controller ranks pending proposals using the
reputations of the lessons they apply or adapt and forms a shortlist
(\texttt{Shortlist} in Algorithm~\ref{alg:search}).
A selector agent chooses one proposal from this shortlist
(\texttt{Select} in Algorithm~\ref{alg:search}).
An implementation agent attempts to implement its strategy
(\texttt{Implement} in Algorithm~\ref{alg:search}).
Each submitted implementation consumes one evaluation attempt, even if
compilation or testing fails.

\begin{figure}[t]
  \centering
  \includegraphics[width=\textwidth]{figures/pipeline/k_lens_pipeline.pdf}
  \caption{Overview of one \klens cycle. Example values from the motivating example cycle in Figure~\ref{fig:motivating-example}.}
  \Description{The K-LENS search loop: search tree, planner agent, rank and shortlist, selector
  agent, implementation agent, Ascend NPU, reflector agent, and measure and update, arranged around
  strategy and coding lessons.}
  \label{fig:overview}
\end{figure}

The controller adds the fastest correct, profiled implementation to the search
tree and compares it with the parent candidate to measure the performance change.
It records a strategy lesson from the parent's conditions, the strategy, and
this result. A reflector agent generates the lesson's summary and extracts coding
lessons from diagnostics and code differences among implementations of the same
proposal (\texttt{Reflect} in Algorithm~\ref{alg:search}). The controller checks the format of the generated
lessons, asking the reflector to retry once if a strategy lesson fails the check, and uses the measured change to update the reputation of the cited
apply/adapt lessons. If no implementation is correct and profiled, the tree
records an all-failed marker and no reputation changes. The agents cannot alter
recorded measurements or assign reputation scores.


\subsection{Notation}
We use lowercase letters for individual candidates, proposals, and lessons, and
uppercase letters for sets. Let $x_i=(f_i,\phi_i,t_i)$ denote a correct, profiled
kernel candidate, where $f_i$ denotes the kernel's source code,
which may span multiple files, $\phi_i$ is its profile fingerprint
(a set of measured profiling features), and $t_i$ is its measured latency.
We write $t(x_i)=t_i$. The search tree $T$ starts from candidate $x_0$ and retains
correct, profiled candidates and markers for cycles in which all attempts failed.

Let $\Pi$ be the set of pending proposals. Each proposal $p\in\Pi$ specifies
a candidate to modify, called its \emph{parent}, and an optimization strategy.
We write $x(p)$ and $m(p)$ for the parent and strategy recorded in $p$.
Let $\mathcal{L}^s$ and $\mathcal{L}^c$ be the sets of strategy and coding
lessons, respectively, and let
$\ell\in\mathcal{L}^s$ denote a strategy lesson.
The set $\mathcal{C}(p)$ contains lessons whose strategies the proposal uses
directly (\emph{apply}) or modifies for the candidate (\emph{adapt}).
$\mathcal{A}(p)$ contains lessons used to rule out an optimization (\emph{avoid}).
The proposal record is $p=(x(p),m(p),\mathcal{C}(p),\mathcal{A}(p))$.
The strategy $m(p)$ specifies, in natural language, the bottleneck it targets
and the profiling measurement that would confirm the expected effect.
The shortlist $H\subseteq\Pi$ contains the proposals considered by the selector.

The sequence $\tau(p)=((a_1,e_1),\ldots,(a_r,e_r))$ records attempts
to implement proposal $p$'s strategy from the same starting candidate,
where $a_j$ is an implementation attempt and $e_j$ is its evaluation result.
Each result records a compilation failure, a correctness failure, or a successful
profile with latency and profiling features. We write $\tau$ when the proposal
is clear from context.
The outcome $y$ is the cycle's fastest correct, profiled implementation,
or an all-failed marker $\bot$ if none succeeds.

Each cycle evaluates at most $R=10$ implementations. Counters $k_1$ and $k_2$
count consecutive evaluations that fail to improve on the parent and on the
cycle's best correct implementation, respectively; compilation and correctness
failures increment both. The cycle ends when a counter reaches $K=5$, after $R$
evaluations, or when the budget is exhausted. The total counter $b$ increments
per submitted implementation, and the search ends when $b$ reaches the budget
$B$.

\subsection{Strategy Lessons and Proposal Selection}
\label{sec:lesson-ranking}
\label{sec:proposal-ranking}

A strategy lesson $\ell$ records the conditions of the candidate being optimized,
the strategy attempted, and the outcome. The planner uses these conditions to
assess the lesson's relevance to another candidate. Table~\ref{tab:strategy-lesson}
lists the lesson fields. The controller computes them deterministically from
proposals, attempts, and measurements, except the summary $z$, generated by
the reflector agent.

\begin{table}[htbp]
  \caption{Fields of a strategy lesson $\ell$.}
  \label{tab:strategy-lesson}
  \label{eq:lesson}
  \small
  \renewcommand{\arraystretch}{1.0}
  \begin{tabular}{@{}p{0.25\linewidth}p{\dimexpr0.75\linewidth-2\tabcolsep\relax}@{}}
    \toprule
    Field & Stored information \\
    \midrule
    Conditions ($q$) & Parent task, bottleneck, profile ratios, and prior strategies. \\
    Strategy ($m$) & The optimization strategy attempted. \\
    Outcome ($o$) & Measurements and change from parent, or no measured result. \\
    Reputation ($\rho$) & Cumulative signed performance credit. \\
    Measurement count ($n$) & Number of performance measurements contributing to $\rho$. \\
    Difficulty ($d$, $w$) & Score summarizing implementation effort, and its accumulated weight. \\
    Source records ($\pi$) & IDs of the candidate, proposal, and attempts that produced the lesson. \\
    Summary ($z$) & A natural-language description of the lesson. \\
    \bottomrule
  \end{tabular}
\end{table}


The planner inspects a summary of the candidate search tree that emphasizes
the best-performing candidates and lists their latencies, profiling features,
and pending proposals. The planner selects a candidate
based on its current performance and opportunities for further improvement,
inspecting source code and full profiles when needed. For the selected candidate,
the controller then presents relevant lesson summaries, ranked by
Gower distance over the recorded candidate context (e.g., bottleneck and prior strategies) and profiling features, the latter compared by
cosine distance. These lessons may come from other candidates in the search or,
with a warm start, from other kernels (Section~\ref{sec:transfer}). The planner
uses these lessons to generate proposals for the selected candidate.

Multiple planning calls can add proposals for different
candidates within a cycle, with at most three pending proposals per candidate.
Candidates at this limit are excluded from planning until a proposal is executed
or retired. Before each selection, the controller requests a proposal for the fastest
candidate if it has none.

\begin{wrapfigure}[22]{R}{0.49\textwidth}
\begin{minipage}{\linewidth}
\scriptsize
\raggedright
\setlength{\algomargin}{1em}
\SetAlFnt{\scriptsize}
\SetAlCapFnt{\scriptsize}
\SetAlCapNameFnt{\scriptsize}
\begin{algorithm}[H]
\caption{\klens search.}
\label{alg:search}
\SetAlgoNoEnd
\SetAlgoNlRelativeSize{-1}
\SetInd{0.3em}{0.5em}
\SetKwInput{KwIn}{Input}
\SetKwInput{KwOut}{Output}
\SetKwFunction{Plan}{Plan}
\SetKwFunction{Shortlist}{Shortlist}
\SetKwFunction{SelectProposal}{Select}
\SetKwFunction{Implement}{Implement}
\SetKwFunction{Reflect}{Reflect}
\KwIn{Correct, profiled $x_0$; budget $B$;
limits $R,K$; initial $\mathcal{L}^s,\mathcal{L}^c$.}
\KwOut{Fastest correct candidate and updated lessons.}
$T\gets\{x_0\}$; $\Pi\gets\emptyset$; $b\gets0$\;
\While{$b<B$}{
  $\Pi\gets\Plan(T,\Pi,\mathcal{L}^s)$\tcp*{planner; \S\ref{sec:lesson-ranking}}
  Recompute all $S_L(p)$, $p\in\Pi$ (Eq.~\ref{eq:proposal-score})\;
  $H\gets\Shortlist(T,\Pi,S_L)$\tcp*{controller; \S\ref{sec:proposal-ranking}}
  $p\gets\SelectProposal(H,T)$\tcp*{selector; \S\ref{sec:proposal-ranking}}
  $x\gets x(p)$; $r\gets\min(R,B-b)$\;
  $(\tau,y)\gets\Implement(p,\mathcal{L}^c,r,K)$\tcp*{\S\ref{sec:attribution}}
  $b\gets b+|\tau|$\tcp*{failures also count}
  Record $\ell'$ from $x,m(p),\tau$, including outcome and source IDs (Table~\ref{tab:strategy-lesson})\;
  $(z,G)\gets\Reflect(\ell',\tau)$\tcp*{reflector; $z$ retried once on format failure}
  \tcp{$z$: strategy summary (\S\ref{sec:lesson-ranking})}
  \tcp{$G$: coding-lesson drafts (\S\ref{sec:coding-lessons})}
  Update $\ell\in\mathcal{C}(p)$ by Eqs.~\ref{eq:credit-update}--\ref{eq:difficulty-update};
  if $y=\bot$, only difficulty changes\;
  \If{$z$ passes the format check}{
    $(\rho_{\ell'},n_{\ell'})\gets(0,0)$\;
    $(d_{\ell'},w_{\ell'},z_{\ell'})\gets(d(\tau),1,z)$\;
    \If{$y\ne\bot$}{
      $(\rho_{\ell'},n_{\ell'})\gets(\Delta(p),1)$\tcp*{Eq.~\ref{eq:delta}}
    }
    $\mathcal{L}^s\gets\mathcal{L}^s\cup\{\ell'\}$\;
  }
  Add coding lessons in $G$ that pass the format check to $\mathcal{L}^c$\;
  Add $y$ to $T$ as a child of $x$, or an all-failed marker if $y=\bot$\;
  $\Pi\gets\Pi\setminus\{p\}$; save search state\;
}
\Return fastest correct candidate in $T$, $\mathcal{L}^s$, $\mathcal{L}^c$\;
\end{algorithm}
\end{minipage}
\end{wrapfigure}

Each proposal cites the presented lessons that informed it as apply, adapt, or
avoid. The controller rejects citations to lessons that were not
presented.
A proposal may cite no lesson when its strategy comes from the source code,
profile, or the model's prior knowledge. The planner records why no lesson was
used; this explanation does not affect ranking.

The controller ranks pending proposals by their \emph{lesson-support score},
computed from the reputation scores of their applied or adapted lessons:
\par\noindent\begin{minipage}{\dimexpr0.51\textwidth-\columnsep\relax}
\begin{equation}
S_L(p)=
\begin{cases}
\displaystyle \frac{1}{|\mathcal{C}(p)|}
\sum_{\substack{\ell\in\mathcal{C}(p)\\n_\ell>0}}\frac{\rho_\ell}{n_\ell},
& |\mathcal{C}(p)|>0,\\[7pt]
0, & \text{otherwise.}
\end{cases}
\label{eq:proposal-score}
\end{equation}
\end{minipage}\par
Here, $|\mathcal{C}(p)|$ denotes the number of lessons cited as apply or
adapt. For a lesson with measurements, $\rho_\ell/n_\ell$ is its average signed
performance credit per measurement. The proposal score averages these
contributions across its applied or adapted lessons. A lesson with no performance
measurements contributes zero. This compares proposals by the average performance
credit of their cited lessons without assigning additional credit for citation
count alone. Avoid citations record why a
direction was ruled out and affect neither ranking nor credit, since the implementation
does not test that direction.

Each proposal stores references to the lessons it cites. After a measured result
updates a lesson's reputation and measurement count, the controller recomputes
the scores of all pending proposals that apply or adapt that lesson. These
proposals may modify different candidates elsewhere in the search tree.
Thus, a result from one candidate can change which proposal is selected for
another, without regenerating or rereviewing the pending proposals.

The controller forms a shortlist from the top-ranked proposals, including one for the fastest
candidate. To keep candidates diverse, it reserves places for proposals for the original candidate
$x_0$ and for candidates not otherwise represented, so lesson scores do not repeatedly favor one
optimization direction. A selector agent chooses one proposal from the shortlist, inspecting the
candidates' profiles and the lesson support.

\subsection{Declare-then-Measure Attribution}
\label{sec:attribution}

Because proposals declare their lesson citations before implementation, measured
credit goes only to lessons declared in advance.
The implementation agent produces implementations for the parent $x(p)$ selected for
optimization, following the strategy in proposal $p$. From the outcome $y$, the fastest correct, profiled implementation, the controller computes the relative latency reduction:
\begin{equation}
  \Delta(p)=\frac{t(x(p))-t(y)}{t(x(p))}.
  \label{eq:delta}
\end{equation}
Thus, $100\Delta(p)\%$ is the percentage reduction in latency relative to the
parent; a negative value means that the implementation is slower.
The controller then updates each cited apply/adapt lesson's reputation $\rho_\ell$,
measurement count $n_\ell$, and implementation difficulty $d_\ell$.
For $c=|\mathcal{C}(p)|>0$ lessons cited as apply or adapt, the update is
\begin{equation}
  (\rho_\ell,n_\ell,d_\ell)\leftarrow
  \begin{cases}
    (\rho_\ell+\Delta(p)/c,\ n_\ell+1,\ \operatorname{Upd}(d_\ell,d(\tau),1/c)),
      & \text{profiled},\\
    (\rho_\ell,\ n_\ell,\ \operatorname{Upd}(d_\ell,5,1/c)),
      & \text{all-failed}.
  \end{cases}
  \label{eq:credit-update}
\end{equation}
The controller divides the measured performance change equally among the $c$
lessons and increments each measurement count by one. The total reputation update
across these lessons is therefore $\Delta(p)$, regardless of citation count.
If no correct, profiled
implementation exists, only difficulty changes; reputation and measurement counts
remain unchanged.

The controller computes the cycle difficulty $d(\tau)$ deterministically, without
an LLM. It ranges from 1 to 4 and increases with the number of attempts needed
to obtain the first correct, profiled implementation; a cycle in which all attempts fail
receives 5. The controller updates the lesson's difficulty
$d_\ell$ by averaging its existing value with the new cycle difficulty $\delta$,
using weight $w$ for the new observation and accumulated weight $w_\ell$:
\begin{equation}
 \operatorname{Upd}(d_\ell,\delta,w)=\frac{w_\ell d_\ell+w\delta}{w_\ell+w},
 \qquad w_\ell\leftarrow w_\ell+w.
 \label{eq:difficulty-update}
\end{equation}
A lesson's initial difficulty has weight 1. For each later proposal, weight 1
is divided equally among its $c$ cited apply/adapt lessons, giving $w=1/c$.
The planner and selector can inspect the difficulty, but it does not
affect the lesson-support score in Equation~\ref{eq:proposal-score}.
A new lesson starts with reputation $\Delta(p)$ and one measurement when profiling
succeeds, or zero reputation and no measurements otherwise. Its initial difficulty
is the cycle difficulty. Avoided and uncited lessons receive no updates. Failed attempts cannot establish whether the strategy
increases or decreases latency. 
% The measured $\Delta(p)$ describes the combined
% effect of the selected implementation's code changes; the update does not estimate
% the contribution of each change or cited lesson separately.

\subsection{Coding Lessons}
\label{sec:coding-lessons}

Coding lessons provide implementation guidance extracted from
the diagnostics and code differences between \klens's attempts to
implement the same proposal. They do not guide planning, proposal
generation, ranking, or selection.
The implementation agent has tools to list coding-lesson
summaries, which are very short and add little to the prompt, and retrieve
specific lessons in full. 

After the cycle, the reflector agent reviews all implementation attempts
  for the same proposal, using code differences
  $\operatorname{diff}(a_u,a_v)$ between attempts, diagnostics,
  and measured results. It extracts lessons that explain how to address
  compilation or testing failures and which conditions code changes must
  preserve, such as synchronization required for data transfers. Each lesson
  states when the guidance applies. These lessons guide implementation in later cycles.

\subsection{Initialization from Other Kernels}
\label{sec:transfer}

We call initialization with strategy lessons from earlier \klens runs on other
kernels a \emph{warm start}. To build a reusable lesson library, we combine
lessons from all runs of each source kernel and keep those from correct,
profiled implementations that reduced latency.
We rank these lessons by parent-relative latency reduction and mean reputation
$\rho_\ell/n_\ell$.
An LLM identifies duplicate or very similar lessons, and we retain the $n_s$
highest-ranked distinct lessons for each kernel. Each lesson retains a reference
to its parent implementation and profile. An LLM
also records which operations or code characteristics each strategy requires,
such as a reduction operation or a particular memory-access pattern. These
descriptions are prepared once and reused across searches.

For each of the 16 kernels in our evaluation, we consider the $n_s$ retained
lessons from each of the other 15 kernels, giving $15n_s$ lessons. The new
search starts from the kernel's original, correct, profiled candidate $x_0$.
An LLM identifies the operations and code characteristics of $x_0$.
The controller deterministically compares these characteristics with each
lesson's recorded requirements and excludes lessons with incompatible
requirements. Then $x_0$ is compared with each remaining lesson's parent candidate $x_j$.
We represent each candidate $x$ by a vector
$\mathbf{u}(x)$ of nine measured pipeline utilization ratios from its fingerprint $\phi$, each in $[0,1]$,
and compute cosine similarity as $s(x_0,x_j)=
\frac{\mathbf{u}(x_0)^\top\mathbf{u}(x_j)}
{\lVert\mathbf{u}(x_0)\rVert_2\,\lVert\mathbf{u}(x_j)\rVert_2}$.
This is the profiling component of the distance used within a search
(Section~\ref{sec:lesson-ranking}); the remaining context, such as prior
strategies, is omitted because it is not comparable across kernels.
The comparison requires complete, nonzero vectors and excludes absolute latency.
The controller ranks the remaining lessons by profile similarity and selects
$n_w$ lessons with distinct optimization mechanisms. Selection is deterministic once the LLM-generated
descriptions are available. In our experiments, $n_s=n_w=5$.

The controller imports the selected lessons before generating proposals for
$x_0$. Imported lessons retain their reputation scores and measurement counts from
earlier runs, guide proposal generation, and contribute to proposal ranking.
% The controller records updates from the new run separately and
% uses the combined values
% $(\rho_\ell,n_\ell)=(\rho_\ell^{\rm past}+\rho_\ell^{\rm new},
% n_\ell^{\rm past}+n_\ell^{\rm new})$ for proposal ranking.
% The labels $\mathrm{past}$ and $\mathrm{new}$ denote contributions from
% the earlier run and the new run, respectively. The new contributions start at
% zero and change according to Section~\ref{sec:attribution}.
The new run also extracts its own strategy and coding lessons, as described in
Sections~\ref{sec:attribution} and~\ref{sec:coding-lessons}.

\section{Implementation}
\label{sec:implementation}

\klens is designed to be reusable across domain-specific languages: the client
manages the agents, search, lessons, and evaluation budget, and a separate backend
manages source files, saved candidate versions, and program execution. The backend
exposes six HTTP/JSON operations: listing, reading, and updating source files;
inspecting and restoring saved versions; and evaluating a candidate through
compilation, testing, and profiling, returning diagnostics and performance
measurements.
This interface separates program editing and evaluation from the search
method, allowing the client's planning, proposal ranking, budget management, and
lesson updates to be reused with other backends. Agents access candidates only
through these operations; only the controller writes measurements, lessons, and
reputation scores. Our \ascendc backend manages source versions and evaluates
candidates through SandboxFusion~\citep{Liu2024FullStackBench} on an Ascend NPU.

The planner, selector, implementation, and reflector agents use the
lightweight Pi coding-agent harness~\citep{Earendil2026Pi}. The controller
interacts with agents only through tool calls and structured outputs, so other
harnesses, such as Codex or Claude Code, can also adapt these
interfaces.
After each cycle, the controller saves the search tree, proposals, lessons,
implementation attempts, measurements, and token usage. These records support
resuming a run and inspecting its recorded decisions.

\section{Evaluation}
\label{sec:evaluation}

We ask four questions:

\begin{description}
  \item[RQ1:] How does \klens compare with the baselines in optimization
  performance under equal evaluation-attempt budgets?
  \item[RQ2:] How does \klens compare with the baselines in optimization
  performance under equal cumulative LLM-token budgets?
  \item[RQ3:] How does initialization with lessons from other kernels affect
  \klens's optimization performance?
  \item[RQ4:] How much does the lesson layer contribute to the agent’s optimization performance?
\end{description}

\subsection{Experimental Setup}

\begin{table}[!t]
  \centering
  \small
  \caption{The \numkernels benchmark operators. SLOC counts editable device sources, including
  header-resident \ascendc code, but excludes build files, tests, generated code, comments, and blanks.}
  \label{tab:benchmark-scale}
  \setlength{\tabcolsep}{3.5pt}
  \footnotesize
  \begin{minipage}[t]{0.485\textwidth}
    \begin{tabular}{llrr}
      \toprule
      Operator & ID & Files & SLOC \\
      \midrule
      \multicolumn{4}{l}{\itshape Architecture-specific operators (12)} \\
      GatherPaKvCache         & GaKV   & 6 & 1{,}326 \\
      InterleaveRope          & IRope  & 5 & 656 \\
      LightningIndexer        & LIdx   & 15 & 4{,}802 \\
      MlaProlog               & MlaP   & 35 & 10{,}382 \\
      MoeFusedTopk            & MoeF   & 2 & 460 \\
      MoeGatingTopK           & MoeG   & 8 & 2{,}838 \\
      MoeTokenUnpermute       & MoeU   & 3 & 437 \\
      QkvRmsNormRopeCache     & QkvC   & 3 & 688 \\
      \bottomrule
    \end{tabular}
  \end{minipage}\hfill
  \begin{minipage}[t]{0.485\textwidth}
    \begin{tabular}{llrr}
      \toprule
      Operator & ID & Files & SLOC \\
      \midrule
      RecurrentGatedDeltaRule & RGDR   & 8 & 1{,}079 \\
      RopeWithSinCosCache     & RopeC  & 8 & 2{,}796 \\
      SparseAttnSharedkv      & SpSKV  & 10 & 5{,}384 \\
      SparseFlashAttention    & SpFA   & 12 & 5{,}669 \\
      \multicolumn{4}{l}{\itshape Common primitives (4)} \\
      AddRmsNorm              & ARMS   & 11 & 3{,}054 \\
      ClippedSwiglu           & CSwG   & 6 & 891 \\
      RmsNorm                 & RMS    & 13 & 3{,}379 \\
      SwiGlu                  & SwiG   & 12 & 919 \\
      \bottomrule
    \end{tabular}
  \end{minipage}
\end{table}

% Draft-only heading spacing; remove for ACM submission.
\vspace{-3pt}
\paragraph*{Environment.}
Candidates are compiled and evaluated in containerized SandboxFusion workers on an
Ascend910 device with CANN 9.1.0. Full environment details are in the artifact.

\vspace{-3pt}
\paragraph*{Model configuration.}
All methods use GPT-5.6 Luna with \texttt{xhigh} reasoning effort.

\vspace{-3pt}
\paragraph*{Benchmark and budget.}
The benchmark contains \numkernels production \ascendc operator packages
(Table~\ref{tab:benchmark-scale}) covering normalization, gated activations, rotary embeddings,
KV-cache movement, latent and sparse attention, mixture-of-experts routing, and gated linear attention.
Their original implementations already contain hardware-specific tuning, including coordinated Cube and Vector stages in fused
attention kernels, making further optimization challenging.

We selected for optimization diversity, not variants of one problem. The benchmark spans
bandwidth-bound elementwise streaming (\textsf{SwiG}, \textsf{RMS}, \textsf{CSwG}, \textsf{ARMS});
irregular gather/scatter (\textsf{GaKV}, \textsf{MoeU}); selection and sorting (\textsf{MoeG},
\textsf{MoeF}); layout transformation fused with cache writes (\textsf{IRope}, \textsf{RopeC},
\textsf{QkvC}); multi-stage attention over tiled workspaces (\textsf{MlaP}, \textsf{SpFA},
\textsf{SpSKV}, \textsf{LIdx}); and a state-carrying sequential recurrence (\textsf{RGDR}). These workloads
have different performance bottlenecks, so an optimization strategy that improves one workload may
not improve another.
We divide the operators into two groups: twelve \emph{architecture-specific operators} for particular
decoder designs and four \emph{common primitives} for normalization and activation. 
% We report results separately for each group.

\begin{wraptable}{r}{0.42\textwidth}
 \centering\footnotesize
 \caption{Main-study configurations.}
 \label{tab:study-size}
 \setlength{\tabcolsep}{3pt}
 \begin{tabular}{lrrr}
 \toprule
 Method & Runs/op. & Eval./run & Runs \\
 \midrule
 \ksearch & 1 & 100 & 16 \\
 \kernelband & 1 & 100 & 16 \\
 \openevolve & 1 & 100 & 16 \\
 \standardcodex & 1 & 100 & 16 \\
 \midrule
 \klens, cold & 3 & 100 & 48 \\
 \klens, warm & 3 & 20 & 48 \\
 \nolessons & 1 & 100 & 16 \\
 \onlycodinglessons & 1 & 100 & 16 \\
 \midrule
 Total & & & \numruns \\
 \bottomrule
 \end{tabular}
\end{wraptable}

Each search starts from the operator's original implementation, which we test and profile before
optimization. All candidates use the same test cases for that operator, covering multiple input shapes and data types. The main comparison allows
\numprofiles evaluation attempts per run. Profiling the original implementation does not count toward
the budget.
Warm-start runs use 20 evaluation attempts to assess how transferred lessons guide initial decisions,
before the search has produced many lessons of its own.
We use \emph{cold} to distinguish standard \klens, which starts without transferred lessons,
from the warm-start version.


For each operator, the artifact records the benchmark input configuration and the measured latency of
the original kernel used as the speedup reference.

Table~\ref{tab:study-size} groups the study configurations.
Because optimization runs require substantial time, hardware use, and LLM tokens, we perform three
runs per operator for \klens and its warm-start version, and one for each baseline
and each ablation (\nolessons and \onlycodinglessons).
The main comparison averages all three cold \klens runs within each kernel.
The same runs provide the cold checkpoints for RQ3 and the variability estimates
in Section~\ref{sec:variability}. The completed GPT-5.6 Luna
study comprises \numruns runs and \numtotalprofiles evaluation attempts, and an estimated 93.57 cumulative run-days.

\paragraph*{Baselines.}
We compare against three specialized kernel optimization approaches and \standardcodex, a
general-purpose coding-agent harness. We retain the specialized baselines' published
search procedures behind a common realization and evaluation interface for \ascendc packages,
replacing hardware-specific profiling features with Ascend measurements. All methods receive
the same task description, editable source files, starting implementation, model and reasoning
settings, and evaluation feedback, and share the same testing and profiling procedures, with
tests and reference data hidden from the agents.

\ksearch is the closest baseline because it also uses a tree-based search
to optimize kernel implementations~\citep{Cao2026K}. We retain its selection rule: among pending proposals whose parent
implementation is available, it selects the proposal with the highest LLM-assigned priority score.
The implementation agent evaluates and repairs implementations of the selected strategy.
We retain its original stagnation setting $K=7$, meaning that a cycle ends after seven consecutive
evaluation attempts fail to improve its best result, or when the total evaluation budget is exhausted.

\kernelband uses a multi-armed bandit to select among six optimization strategies~\citep{Ran2025KernelBand}.
We retain its published strategy set and selection rule, with Ascend measurements replacing
NVIDIA-specific ones. The Ascend profiler lacks the measurements required by its pipelining filter,
so we retain pipelining among the available strategies.

\openevolve (commit \texttt{b28ac15}) evolves a database of candidate programs, initialized with the
measured starting kernel, using its default population settings~\citep{Sharma2025OpenEvolve}.
As with the other methods, each run receives \numprofiles evaluation attempts.

For \kernelband and \openevolve, we replace single-response code generation with the implementation
agent used by \ksearch. The agent can edit code over multiple turns to coordinate changes across the
files in each operator package (Table~\ref{tab:benchmark-scale}). Each method still selects the parent
and optimization strategy, which remain fixed during editing. The agent cannot evaluate candidates
while editing. As in the original \kernelband and \openevolve methods, the backend evaluates each
completed candidate once, without a repair step based on that result.


\standardcodex uses its standard shell and editing tools, with a client
for the shared compile--test--profile backend. Given the same source code, task information,
model, and reasoning level as the other methods, but no kernel-specific search
controller or lesson set, it is prompted to improve performance while preserving
behavior and told its evaluation budget. The artifact records the prompts and the
evaluation and token logs.


\vspace{-3pt}
\paragraph*{Accounting and metrics.}
Every submitted implementation counts toward the evaluation-attempt budget,
including compilation and test failures.
We compare total LLM-token budgets, without separate limits
for input, cached-input, and output tokens.
Cached input dominates token use for four methods, whereas \ksearch uses more
uncached input, so equal total-token budgets may cost more when cached tokens are priced lower.
The artifact reports token composition and usage by method; token use at each
checkpoint is read from cumulative token counts.

For kernel $i$, method $m$, run $r$, and budget $b$, let
$s_{i,m,r,b}=\bar t_i/t^{\mathrm{best}}_{i,m,r,b}$ denote the
best-so-far speedup relative to the original implementation.
Here, $\bar t_i$ is a shared reference per kernel: the mean of its original-implementation
measurements. The artifact lists the measurements included in each comparison.
$t^{\mathrm{best}}_{i,m,r,b}$ is the lowest execution time among the original implementation
(available at zero budget) and correct candidates within budget $b$.
We first average speedups across independent runs for each kernel,
then take the geometric mean (GM) across kernels:
\begin{equation}
 \bar s_{i,m,b}=\frac{1}{R_m}\sum_{r=1}^{R_m}s_{i,m,r,b},\qquad
 \operatorname{GM}_m(b)=\exp\!\left(\frac{1}{N}\sum_{i=1}^{N}\log\bar s_{i,m,b}\right),
 \label{eq:aggregates}
\end{equation}
Here, $N=16$ is the number of kernels and $R_m$ is the number of repetitions per kernel for method $m$:
$R_m=3$ for standard and warm-start \klens, and $R_m=1$ for each baseline,
\nolessons, and \onlycodinglessons.
We use the same calculation for evaluation-attempt and cumulative LLM-token budgets.
For \onlycodinglessons, we report attempt-budget results; final token consumption is in the artifact.
Its normalized records lack candidate-level token counters for token checkpoints.
Each kernel receives equal weight in log-speedup~\citep{Fleming1986HowNot,Ran2025KernelBand}.
The geometric mean summarizes normalized performance across this fixed suite;
it does not estimate end-to-end application speedup.
A speedup factor $s$ corresponds to
$100(s-1)$\% speedup and $100(1-1/s)$\% latency reduction.

Pass measures the fraction of total evaluation attempts that compile and pass testing.
Improved measures the fraction of total evaluation attempts per run that produce correct candidates
faster than the shared baseline. We report per-operator speedups and per-category and overall GMs, with uncertainty across
\klens runs estimated in Section~\ref{sec:variability}.


\subsection{RQ1: Equal Evaluation-Attempt Budgets}
\begin{figure*}[t]
  \centering
  \includegraphics[width=\textwidth]{figures/evaluation/profile-trajectories.pdf}
  \caption{GM speedup by evaluation-attempt budget: latest correct (left)\textsuperscript{*}
  and best-so-far (right) candidates. Shading shows pointwise approximate 95\%
  intervals for \klens (Section~\ref{sec:variability}).}
  \Description{Two step plots compare five methods over 100 evaluation attempts. The left
  chart shows the geometric-mean speedup of the latest correct candidates and the right chart the
  geometric-mean speedup of the best-so-far candidates across all \numkernels operators.}
  \label{fig:profile-trajectories}
\end{figure*}
\begingroup
\renewcommand{\thefootnote}{\fnsymbol{footnote}}
\footnotetext[1]{For readability, the left panel omits an \openevolve \textsf{RopeC} outlier
(attempt 72; about $20\times$ the baseline latency), retaining the preceding correct candidate
through attempt 75. Evaluation counts, tables, and best-so-far results are unchanged.}
\endgroup

\begin{table}[t]
  \centering
  \small
  \caption{Best-so-far GM speedup across \numkernels kernels by evaluation budget;
  Arch/Prim, Improved, and Pass are reported at \numprofiles attempts.
  Grey shading marks final-budget columns.}
  \label{tab:profiles}
  \setlength{\tabcolsep}{3pt}
  \begin{tabular}{lccccc@{\hspace{6pt}}cc@{\hspace{6pt}}cc}
    \toprule
    & \multicolumn{5}{c}{GM speedup at attempt} & \multicolumn{2}{c}{By category} & & \\
    \cmidrule(lr){2-6}\cmidrule(lr){7-8}
    Method & 20 & 40 & 60 & 80 & \tikzmarknode{profiles-budget-top}{\numprofiles} & Arch & Prim & Improved & Pass \\
    \midrule
    \ksearch & 1.0772 & 1.0826 & 1.0842 & 1.0909 & 1.0964 & 1.1149 & 1.0426 & 56.50\% & 86.94\% \\
    \kernelband & 1.0900 & 1.0947 & 1.0970 & 1.1097 & 1.1138 & 1.1286 & 1.0705 & 43.06\% & 79.94\% \\
    \openevolve & 1.0697 & 1.0755 & 1.0936 & 1.0975 & 1.0982 & 1.1102 & 1.0631 & 44.63\% & 85.13\% \\
    \standardcodex & 1.0762 & 1.0953 & 1.1095 & 1.1106 & 1.1110 & 1.1157 & \textbf{1.0971} & 57.88\% & \textbf{92.44\%} \\
    \klens & \textbf{1.1762} & \textbf{1.2004} & \textbf{1.2100} & \textbf{1.2213} & \tikzmarknode{profiles-budget-bottom}{\textbf{1.2240}} & \textbf{1.2718} & 1.0913 & \textbf{74.48\%} & 88.35\% \\
    \bottomrule
  \end{tabular}%
  \begin{tikzpicture}[remember picture,overlay]
    \node[draw=none,fill=black,fill opacity=0.08,inner xsep=3pt,inner ysep=3pt,
      fit=(profiles-budget-top)(profiles-budget-bottom)] {};
  \end{tikzpicture}
\end{table}

\klens achieves the highest GM speedup at all five checkpoints in Table~\ref{tab:profiles},
with about twice the percentage speedup of the strongest baseline throughout
(+22.40\% versus +11.38\% for \kernelband at \numprofiles attempts). The lead is
established early: at 20 attempts \klens reaches +17.62\%, more than any baseline
achieves at 100.

Figure~\ref{fig:profile-trajectories} shows frequent declines in latest-correct
performance for \kernelband and \openevolve, with only intermittent gains
in their best-so-far results. They also have the lowest Pass (79.94\%, 85.13\%)
and Improved (43.06\%, 44.63\%) rates. Because Improved counts attempts that
beat the original, dividing it by Pass gives the share of passing attempts that
improve: 54\% and 52\%, against 84\% for \klens. We interpret this as a
difference in strategy quality: without measured evidence on which strategies work,
these methods try riskier, more diverse rewrites whose correct candidates often do not
improve speed, whereas \klens favors strategies more likely to improve.

\ksearch's Pass rate (86.94\%) is close to \klens's (88.35\%), yet its Improved rate is
56.50\% versus 74.48\%: the gap lies in which proposals are implemented, not in whether
attempts succeed. \ksearch shares \klens's tree and plan/implement split but ranks by
LLM-assigned priority, and trails by 12.8 points in GM speedup (+9.64\% versus +22.40\%),
which points to the value of ranking by measured lesson evidence.

\klens performs best on 12 of the 16 kernels, while \standardcodex performs best on three
(\textsf{RGDR}, \textsf{RMS}, and \textsf{SwiG}; Figure~\ref{fig:per-kernel-bars}) and
\kernelband on one (\textsf{GaKV}, by 0.5 points). \standardcodex leads
on common primitives (+9.71\%, compared with +9.13\% for \klens).
\standardcodex is effective at producing valid implementations, achieving
the highest Pass rate of 92.44\%. Its lower Improved rate
(57.88\%, compared with 74.48\% for \klens) shows that implementation
success alone does not ensure improvement: only 63\% of its passing attempts
beat the original, against 84\% for \klens. In our baseline,
long-term planning and reuse of previous results depend on the agent's
use of its session history. \klens provides explicit rules for storing
measured optimization outcomes and using them to rank subsequent proposals.

The simpler structure of common primitives may allow \standardcodex to find effective
optimizations through direct edits and evaluation feedback. Their widespread use may also mean
more training examples of these computations and optimization patterns.
\klens's advantage is larger on architecture-specific operators: +27.18\%, compared with
+12.86\% for the strongest baseline, consistent with lesson reuse paying off most when
several optimization decisions interact.

\begin{figure*}[t]
  \centering
  \includegraphics[width=\textwidth]{figures/evaluation/per-kernel-bars.pdf}
  \caption{Per-kernel best-so-far speedup and overall GM after \numprofiles evaluation attempts.
  IDs from Table~\ref{tab:benchmark-scale}. Grey shading marks overall GM results in figures.}
  \Description{A grouped bar chart compares five methods across 16 kernels: twelve
  architecture-specific operators and four common primitives, followed by a shaded geometric-mean
  group. K-LENS has the highest speedup on twelve kernels, Codex on three, and KernelBand on one.}
  \label{fig:per-kernel-bars}
\end{figure*}

\begin{klensresultbox}
\textbf{RQ1:} At \numprofiles attempts, \klens achieves approximately twice the percentage
speedup of the second-best method, \kernelband: +22.40\% versus +11.38\%.
This approximately twofold advantage holds at all five reported checkpoints.
\end{klensresultbox}

\subsection{RQ2: Equal LLM-Token Budgets}

% \begin{figure}[t]
%   \centering
%   \includegraphics[width=\textwidth]{figures/evaluation/token-trajectories.pdf}
%   \caption{GM speedup across \numkernels kernels by cumulative LLM-token budget: latest correct
%   candidate (left) and best-so-far candidate (right). Shading shows pointwise approximate 95\%
%   intervals for \klens (Section~\ref{sec:variability}); baselines use one run per kernel.
%   Through 50M tokens, the left panel samples \kernelband and \openevolve every 0.5M tokens,
%   omitting intermediate fluctuations. Other curves and the numerical analysis use all recorded events.}
%   \Description{Two step plots compare the latest-correct and best-so-far geometric-mean speedups
%   of five methods across 16 kernels under cumulative budgets up to 100 million LLM tokens.}
%   \label{fig:token-trajectories}
% \end{figure}

\begin{table}[t]
  \centering
  \small
  \caption{Best-so-far GM speedup by cumulative LLM-token budget. M denotes million tokens.}
  \label{tab:tokens}
  \begin{tabular}{lccccc}
    \toprule
    Method & 10M & 20M & 30M & 50M & \tikzmarknode{tokens-budget-top}{100M} \\
    \midrule
    \ksearch & 1.0750 & 1.0770 & 1.0814 & 1.0841 & 1.0898 \\
    \kernelband & 1.0824 & 1.0863 & 1.0984 & 1.1084 & 1.1138 \\
    \openevolve & 1.0714 & 1.0923 & 1.0963 & 1.0970 & 1.0977 \\
    \standardcodex & 1.0199 & 1.0586 & 1.0696 & 1.0838 & 1.0959 \\
    \klens & \textbf{1.1244} & \textbf{1.1595} & \textbf{1.1733} & \textbf{1.1960} & \tikzmarknode{tokens-budget-bottom}{\textbf{1.2177}} \\
    \bottomrule
  \end{tabular}%
  \begin{tikzpicture}[remember picture,overlay]
    \node[draw=none,fill=black,fill opacity=0.08,inner xsep=3pt,inner ysep=3pt,
      fit=(tokens-budget-top)(tokens-budget-bottom)] {};
  \end{tikzpicture}
\end{table}

We apply common token caps retrospectively to the recorded runs, reporting the best correct candidate available within each cap.
Wall-clock time is not matched because evaluation duration also depends on failures, kernel speed,
queues, and network conditions. We compare the fastest correct kernel each method produces
within each token budget.


% Figure~\ref{fig:token-trajectories} shows optimization progress under cumulative token budgets.
Table~\ref{tab:tokens} shows that \klens leads at all five token budgets.
\kernelband is closest at 10M tokens, where \klens's percentage gain is
1.51 times as large. This advantage increases to 1.91 times at 100M
tokens (+21.77\% versus +11.38\%). Thus, the gap widens as the token
allowance increases.

At 10M tokens, \klens (+12.44\%) already exceeds every baseline at 100M tokens
(at most +11.38\%). Baselines flatten: \openevolve, the strongest at 20M, improves
only from +9.23\% to +9.77\% thereafter, while \klens continues to gain
(+15.95\% at 20M, +21.77\% at 100M). Additional tokens are converted into
improvement mainly by \klens, whose ranking needs no LLM calls.

\standardcodex has the lowest speedup at 10M tokens (+1.99\%), despite
its high Pass rate in RQ1. Early in the search, code inspection and
checks between evaluations may consume tokens without measured
improvements. Its percentage
speedup rises to +9.59\% at 100M tokens, still below \kernelband.
This shows why optimization under a token budget required separate evaluation.

\begin{klensresultbox}
\textbf{RQ2:} \klens leads at every reported token budget. At 100M tokens,
its percentage gain is 1.91 times that of the strongest baseline
(+21.77\% versus +11.38\%). At 10M, it already exceeds all baselines at 100M.
\end{klensresultbox}

\subsection{RQ3: Effect of Cross-Kernel Lesson Transfer}
\label{sec:rq3}

We compare three warm runs per kernel with the three cold
runs from RQ1 at 10 and 20 evaluation attempts. Warm runs import five
lessons from other kernels before optimization
(Section~\ref{sec:transfer}); both conditions then use the same search
procedure. Cold GMs in this section differ slightly from Table~\ref{tab:profiles}
because the shared baseline also includes warm-run baseline measurements.

Warm start achieves higher mean speedup on 15 of 16 kernels at both
checkpoints (Figure~\ref{fig:transfer-effects}). At 10 attempts,
GM speedup is 1.2230 versus 1.1557 for cold search; at 20 attempts,
it is 1.2845 versus 1.1758 (Table~\ref{tab:transfer}). These correspond
to 1.43 and 1.62 times the cold percentage gain, respectively.

\begin{table}[t]
  \centering
  \footnotesize
  \setlength{\tabcolsep}{3pt}
  \begin{minipage}[t]{0.36\textwidth}
  \centering
  \caption{Warm-start GM speedup.}
  \label{tab:transfer}
  \begin{tabular}{@{}lcc@{}}
    \toprule
    Metric & Cold & Warm \\
    \midrule
    GM @10 & 1.1557 & 1.2230 \\
    GM @20 & 1.1758 & 1.2845 \\
    Arch-specific @20 & 1.2161 & 1.3139 \\
    Primitives @20 & 1.0626 & 1.2003 \\
    \bottomrule
  \end{tabular}
  \end{minipage}\hfill
  \begin{minipage}[t]{0.62\textwidth}
  \centering
  \caption{Lesson-layer ablation at \numprofiles attempts (RQ4).}
  \label{tab:ablation}
  \begin{tabular}{@{}lccc@{}}
    \toprule
    Metric & \klens & \nolessons & \onlycodinglessons \\
    \midrule
    GM speedup & 1.2240 & 1.1315 & 1.1452 \\
    Improved & 74.48\% & 64.88\% & 72.44\% \\
    Attempts to first pass & 1.31 & 2.06 & 1.19 \\
    Compile-failure rate & 5.38\% & 2.94\% & 2.81\% \\
    \bottomrule
  \end{tabular}
  \end{minipage}
  \par\smallskip
  \begin{minipage}{\textwidth}\footnotesize
  Table~\ref{tab:ablation} uses a common fixed baseline per kernel.
  Compilation failures are 258/4,800 for \klens, 47/1,600 for
  \nolessons, and 45/1,600 for \onlycodinglessons. The two ablations
  have one run per kernel; \klens has three.
  \end{minipage}
\end{table}

\begin{figure*}[t]
  \centering
  \includegraphics[width=\textwidth]{figures/warm_start/rq3_transfer_effects.pdf}
  \caption{Warm-minus-cold percentage speedup at 10 and 20 attempts,
  in percentage points. Per-kernel values compare mean speedups;
  Overall compares GMs. Positive values favor warm start.}
  \Description{Two panels show warm-minus-cold differences for
  16 kernels and the overall GM. Warm start has higher mean speedup
  on 15 kernels at both checkpoints.}
  \label{fig:transfer-effects}
\end{figure*}

Warm start raises the percentage gain at 20 attempts from +17.58\% to +28.45\%, by 10.88 points.
Excluding any single kernel leaves 7.98--11.78 points, so no one kernel accounts for the effect.
The per-category gains in Table~\ref{tab:transfer} are less stable: the gain on common primitives
rests mainly on \textsf{RMS}.

At 20 attempts, warm start also exceeds cold search at 100 attempts on 10 of 16
kernels (GM 1.2845 versus 1.2235, 4.99\% higher with one fifth of the budget).
Transferred lessons thus supply measured strategies that the planner can adapt before
local results accumulate, and that a cold search does not otherwise find within a
larger budget.

\begin{klensresultbox}
\textbf{RQ3:} Warm runs yield 1.43 and 1.62 times the cold percentage
gain at 10 and 20 attempts. Warm at 20 outperforms cold at 100 on 10 of 16 kernels.
\end{klensresultbox}

\subsection{RQ4: Contribution of the Lesson Layer}
\label{sec:rq4}

\nolessons uses an LLM to select parent--strategy proposals without
lessons and therefore without lesson-based ranking. It retains
\klens's implementation agent, evaluator, and budget, but removes reflection.
We also evaluate \emph{\onlycodinglessons}, which is identical to \nolessons
except that it extracts and uses coding lessons. Given the evaluation cost,
we compare full \klens with two ablations, \nolessons and \onlycodinglessons,
to assess the strategy-lesson mechanism, lessons together with the ranking they drive,
against coding guidance alone.

At \numprofiles attempts, \klens achieves a 70\% larger percentage speedup than
\nolessons (+22.40\% versus +13.15\%), with higher mean speedup on 14 of 16 kernels
(Figure~\ref{fig:lesson-effects}). Without lessons, the percentage gain is closer
to the other baselines' +9.64--11.38\% (Table~\ref{tab:profiles}) than to \klens's
+22.40\%.
Both ablations retain \klens's tree, evaluator, and budget. \nolessons removes all
lessons; \onlycodinglessons retains coding lessons but removes strategy lessons and their ranking. Adding coding guidance alone yields a modest gain, from +13.15\% to +14.52\%.

Compared with \nolessons, \klens has a higher compilation-failure rate (5.38\% versus
2.94\%) but a higher Improved rate (74.48\% versus 64.88\%), and still reaches a
first passing implementation in fewer attempts (1.31 versus 2.06; Table~\ref{tab:ablation}).
Combining earlier strategies with existing optimizations requires coordinated
multi-file edits, riskier but yielding faster kernels; we interpret the
short time to first pass as coding lessons repairing these failures quickly. Excluding any one kernel, \klens still achieves 5.67--8.89\% higher GM
speedup, so no single kernel accounts for the advantage.


\begin{figure}[t]
  \centering
  \includegraphics[width=\textwidth]{figures/evaluation/lesson-effects-horizontal.pdf}
  \caption{\klens speedup advantage (\%) over \nolessons (blue) and
  \onlycodinglessons (green) at \numprofiles attempts. }
  % Per-kernel points
  % compare mean speedup across three \klens runs with each ablation's single
  % run; the final pair compares GMs.}
  \Description{Per-kernel speedup advantages of K-LENS over No-Lessons
  and Only-Coding. }
  \label{fig:lesson-effects}
\end{figure}

% At 100 attempts, using the lowest speedup from the three \klens runs
% for each kernel still yields a 4.39\% higher GM speedup than \nolessons.

Adding coding lessons reduces mean attempts to first pass from 2.06 to 1.19 and raises the
Pass rate from 85.25\% to 91.81\%. Of \klens's 9.25-point gain in percentage speedup over
\nolessons, 1.37 points (15\%) are attributable to coding lessons and 7.88 (85\%) to strategy
lessons with lesson-based ranking; this decomposes observed differences and does not separate
components that may interact. \onlycodinglessons keeps \klens's tree,
agents, and coding guidance yet stays closer to the baselines (+9.64--11.38\%) than to \klens,
so the harness is not the source of the gain. In our recorded searches, scores that incorporate
later measurements align better with realized parent-relative improvement than the scores lessons
carried when created.

\begin{klensresultbox}
\textbf{RQ4:} At \numprofiles attempts, \klens achieves 1.70 times the percentage
speedup of \nolessons (+22.40\% versus +13.15\%), supporting the lesson layer's
contribution to optimization. Coding lessons mainly improve implementation reliability, with only a
  modest contribution to overall optimization gains.
\end{klensresultbox}

\subsection{Variation across Independent Runs}
\label{sec:variability}

Three independent \klens runs per kernel give mean speedup $\bar s_i$ and sample variance
$\hat\sigma_i^2$. We summarize run-to-run variation by the median within-kernel variance,
the median coefficient of variation, and an approximate 95\% confidence interval for the suite GM
(Equation~\ref{eq:aggregates}). With
$v_i=\hat\sigma_i^2/(3\bar s_i^2)$ and $\hat\mu=N^{-1}\sum_i\log\bar s_i$,

\begin{equation}
 \widehat{\operatorname{SE}}(\hat\mu)=\frac{\sqrt{\sum_i v_i}}{N},\qquad
 \mathrm{CI}_{95}=\exp\!\left(\hat\mu\pm t_{0.975,\nu}\,
          \widehat{\operatorname{SE}}(\hat\mu)\right),
 \label{eq:llm-variability}
\end{equation}
where $\nu=(\sum_i v_i)^2/\sum_i(v_i^2/2)$ is the Welch--Satterthwaite effective degrees of
freedom~\citep{NIST1994Uncertainty}. The approximation assumes independent searches and modest
within-kernel variation. The interval characterizes \klens's run-to-run variability
on the fixed benchmark suite.

At \numprofiles evaluation attempts, the median within-kernel sample variance of speedup is
$7.052\times10^{-4}$, with a median coefficient of variation of $2.46\%$.
The suite GM speedup is $1.2240\times$, with an approximate 95\% confidence interval of
$[1.2037,\,1.2447]$ using Equation~\ref{eq:llm-variability}.
The interval extends about 1.7\% on either side of the GM estimate, whereas \klens's GM speedup
is 9.90\% higher than that of the strongest baseline, \kernelband ($1.1138\times$).
The interval describes \klens only: one \klens run has a standard error of about 1.3\% (log scale),
so the gap to \kernelband exceeds three standard errors even if baseline runs were twice as variable,
and each of the three \klens runs alone beats every baseline.
The same analysis for the warm-start runs at 20 attempts gives a median
within-kernel variance of $2.684\times10^{-4}$ and a 95\% suite-GM interval of $[1.2525,\,1.3174]$.

\subsection{Generalization to Another Model}
\label{sec:model-generalization}

\begingroup
\setlength{\intextsep}{3pt}
\setlength{\columnsep}{7pt}
\begin{wraptable}[9]{l}{0.42\textwidth}
  \centering\footnotesize
  \captionsetup{font=footnotesize,skip=2pt}
  \caption{GM speedup on six kernels with Kimi K3.}
  \label{tab:model-generalization}
  \setlength{\tabcolsep}{3pt}
  \begin{tabular}{@{}lcc@{}}
    \toprule
    Method & \shortstack{100\\attempts} & \shortstack{50M\\tokens} \\
    \midrule
    \ksearch       & 1.1231 & 1.1231 \\
    \kernelband    & 1.0724 & 1.0724 \\
    \openevolve    & 1.0645 & 1.0645 \\
    \standardcodex & 1.1972 & 1.1971 \\
    \klens        & \textbf{1.5033} & \textbf{1.3909} \\
    \bottomrule
  \end{tabular}
\end{wraptable}
We repeat the comparison with Kimi K3 using \texttt{low} reasoning.
To target the most complex kernels, we set a 3{,}000-SLOC threshold, which leaves six:
\textsf{SpSKV}, \textsf{LIdx}, \textsf{MlaP}, \textsf{RMS}, \textsf{SpFA}, and
\textsf{ARMS} (Table~\ref{tab:benchmark-scale}). Evaluation cost limits the study to
one run per method and kernel.
Table~\ref{tab:model-generalization} reports best-so-far GM speedup at
\numprofiles evaluation attempts and 50M tokens for the same six kernels.
With low reasoning, many runs finish below 100M
tokens; we therefore report the token-budget comparison at 50M.

At \numprofiles evaluation attempts, \klens achieves
2.55 times the percentage speedup of the strongest baseline, \standardcodex
(+50.33\% versus +19.72\%). At 50M tokens, the ratio is 1.98
(+39.09\% versus +19.71\%).
The strongest baseline changes with the model (\standardcodex here), yet \klens stays
first under both budgets in these independent runs. These results provide evidence that
the advantage extends to a second model family on the evaluated kernels.
% although one run per method and kernel does not establish
% run-to-run variability.
\par
\endgroup
\section{Discussion and Threats to Validity}
\label{sec:threats}

\emph{External validity.} Evaluation cost limits model and benchmark coverage:
the 192 Luna runs consumed tokens on the order of tens of billions and an
estimated 93.57 cumulative run-days. We selected GPT-5.6 Luna to balance coding
capability and inference cost. A comparison with Kimi K3 on six complex kernels exceeding
3{,}000 SLOC (Section~\ref{sec:model-generalization}) also favors \klens, providing
evidence beyond the primary model family. Transfer may be easier between closely
related kernels. We therefore selected 16 kernels covering distinct computations,
implementation sizes, and bottlenecks, including architecture-specific operators
and common primitives, and report results for all of them.
% TODO: state selection criteria and that they were fixed before the experiments.
% \textcolor{red}{[Dr. Meng: Please add more details about benchmark selection.]}
Nevertheless, their shared library, language, and hardware limit generalization. Exposure to
these kernels in training data yields no speedup, since every method starts from
them.

\emph{Internal validity.} Random model outputs and timing noise can affect the
comparison. Three independent \klens runs per kernel show modest observed
variability, with a median coefficient of variation of 2.46\%
(Section~\ref{sec:variability}). Baselines and ablations use one run per kernel, so their variability
is not characterized.
% TODO: if available, add per-kernel wins, e.g. "\klens leads on X of 16 kernels".
Timing noise is small relative to the gaps: candidates are profiled in containerized
sandboxes, each NPU serves only one run at a time, and every profile discards the first three
measurements as warm-up and averages the rest; % TODO: confirm warm-up/averaging is what the code does.
the spread of repeated baseline measurements shifts the suite GM by about 0.7\%.
All methods use the same starting implementations, hidden
tests, evaluation backend, and budgets, and agents cannot inspect the correctness tests. After the searches were over, we re-ran the best
candidate they found per kernel against the original kernel in a separate session, outside the
budget, and the improvements they reported were reproduced. For \ksearch, \openevolve, and \kernelband, we
adapt the code-editing and evaluation interfaces to \ascendc and release these
adaptations. Multi-turn editing gives \kernelband and \openevolve additional
implementation support compared with their original single-response generation.
\kernelband's pipelining filter requires counters unavailable from the Ascend
profiler, so we retain pipelining among its available strategies rather than
filtering it out.
\klens's parameters ($R$, $K$, shortlist size, $n_s$, and $n_w$) were fixed
before the evaluation rather than tuned per kernel. % TODO: confirm.

\emph{Design choices.} Dividing measured performance changes equally among cited
lessons is a design choice because the measured outcome does not identify each
lesson's contribution. Because the measurement count still increments by one,
co-citation can lower a lesson's mean credit. Fewer than 10\% of evaluated proposals cite multiple
lessons, limiting how often this choice affects attribution. The lesson ablation
removes lessons and their ranking together; \ksearch, which ranks by LLM
priority, is the alternative.
Building the warm-start bank is excluded from each target's budget, amortized
across targets, and needs no kernel evaluations or device time, only an LLM pass of
under 1M tokens over 80 lessons the source searches already produced. Equal token budgets do not imply equal costs: cached input dominates
token use for four methods, whereas \ksearch uses more uncached
input, which can cost more when cached tokens are cheaper. Given the
computational cost, future work includes broader evaluation across models,
reasoning settings, libraries, and hardware, and additional repetitions for the
baselines and ablations.


\section{Related Work}

\noindent\textbf{Measured kernel search.}
Like autotuners, \klens uses measurements to guide
search~\citep{Ansel2014OpenTuner,Chen2018Learning,Zheng2020Ansor,Adams2019Halide},
but ranks natural-language optimization proposals rather than points in a fixed
search space.
\ksearch separates optimization planning from implementation in a search tree and
uses an LLM to revise pending-action priorities~\citep{Cao2026K}.
\kernelband uses hardware-aware bandits to select among a fixed set of
optimization strategies, sharing reward statistics among candidates with similar
profiling traces~\citep{Ran2025KernelBand}. ReGraphT instead builds a graph of
optimization techniques from LLM-generated trajectories and backpropagates
validation and performance reward through it during search~\citep{Gong2025From}.
In both, credit attaches to strategies or techniques in a structure fixed or
built before the search, not to lessons extracted during it. Population and multi-agent systems
preserve programs, diagnostics, or strategy
histories~\citep{Chen2025cuPilot,Sun2026KernelSkill,Gai2026Optimizing,Novikov2025AlphaEvolve,Sharma2025OpenEvolve,Liao2025KernelEvolve};
profiler-guided and evolutionary systems drive revision with
measurements~\citep{Lei2025PRAGMA,Guo2025EvoEngineer,Du2026Kernel}.
\klens instead credits lessons extracted during the search through the proposals
that cite them. The controller uses parent-relative measurements to update applied
or adapted lessons and recomputes the scores of pending proposals that cite them,
including those for other candidates, without regenerating proposals or
LLM-based rescoring.

\noindent\textbf{Experience for optimization agents.}
LessonL updates lesson effectiveness using later speedups and selects guidance
for code generation using effectiveness and similarity to the original
code~\citep{Liu2025Lessons}; \klens instead uses reputation to rank which
proposal is implemented, crediting lessons cited before implementation.
EvoKernel updates memory using a refinement reward that compares latency with the
best-so-far implementation and penalizes infeasibility~\citep{Zheng2026Cold}.
KernelBlaster selects state-conditioned CUDA optimizations using stored scores
but revises its strategy knowledge through LLM-based performance
analysis~\citep{Dong2026KernelBlaster}.
daVinci-kernel selects and summarizes execution-validated
skills~\citep{Fu2026daVinciKernel}; KernelSkill joins long-term knowledge with
task-local memory~\citep{Sun2026KernelSkill}; and AccelOpt generalizes selected
slow--fast kernel pairs into reusable plans~\citep{zhang2026accelopt}. For
transfer, \klens selects lessons by requirement compatibility and profile
similarity and continues updating their reputations on the new kernel.
Scored experience is thus not new in itself; \klens's distinction is
declare-then-measure credit: lessons cited before implementation receive
parent-relative credit, which reranks pending proposals for other candidates
without LLM rescoring.

\noindent\textbf{Agent memory.}
Reflexion retains outcome-derived verbal feedback for later
decisions~\citep{Shinn2023Reflexion}, while ExpeL extracts experience
that can transfer across tasks~\citep{Zhao2023ExpeL}. For code optimization,
SBLLM combines retrieved optimization patterns with execution-based
candidate selection~\citep{Gao2025SBLLM}. Repeated language-model rewriting
can degrade useful memory~\citep{Zhang2026Useful}; \klens does not rewrite
recorded lessons and updates only their numerical scores.

\noindent\textbf{Ascend systems.}
AscendOptimizer identifies optimizations whose removal increases latency, then
combines the resulting patterns with profiling-guided host-side
search~\citep{Wu2026AscendOptimizer}. Hawk maintains structured AscendC knowledge
online, refining its applicability or removing entries after
failures~\citep{Wen2026Hawk}; when no implementation is measured, \klens leaves
reputation unchanged and updates only implementation difficulty.
AscendKernelGen, AscendCraft, and AgenticCANN strengthen model knowledge,
structural validity, or knowledge
orchestration~\citep{Cao2026AscendKernelGen,Wen2026AscendCraft,Qiu2026AgenticCANN};
\klens instead accumulates measured knowledge during search and is complementary
to these approaches. NPU benchmarks report challenges in kernel correctness and
performance~\citep{Wen2025MultiKernelBench,Gao2026CANN}.


\section{Conclusion}

\klens treats every measurement as something to learn from. It records what each
optimization did to a kernel as a lesson, uses these lessons to propose new
optimizations, and ranks proposals by the measured gains of the lessons they cite.
Each new result updates the scores of the lessons it relied on and reranks the
pending proposals, without regenerating them. On \numkernels \ascendc kernels,
\klens nearly doubles the strongest baseline's gain at \numprofiles evaluation
attempts (22.40\% versus 11.38\%) and leads at every token budget. Without the
lesson layer, the gain falls to 13.15\%, so lessons account for most of its
advantage over the baselines. Starting a search with lessons from other kernels
raises the gain at 20 attempts from 17.58\% to 28.45\%.

\section*{Data Availability}

The anonymized replication package is available at \url{https://figshare.com/s/7000a3edb4be81c6ea22}. It contains the 16-kernel benchmark, configurations and prompts, client code, proposal and lesson records, evaluation logs,
token accounting, analysis scripts and an experimental details document with the full environment, baseline, and accounting
configuration. The best optimized source version is included for each of the 16 operators,
with its recorded baseline and optimized performance.

\bibliographystyle{ACM-Reference-Format}
\bibliography{references}

\end{document}
